% Argument revision after adversarial review, 2026-09-05. Hand-edited.
% Sources and claim boundaries: positioning.md and reviews/2026-09-05-adversarial/.
% Empirical result preview remains deferred at the user's request.
\section{Introduction}
\label{sec:introduction}

Activation sparsity can reduce language-model inference cost when specialized
kernels skip products and weight accesses associated with zero activations.
Recent work shows that useful sparsity can be exposed in pretrained models
through thresholding or encouraged during training, with measured acceleration
in both cases~\citep{liu2024teal,cetin2026sparser}. These successes create a
practical design problem: how should a researcher combine and place
sparsification interventions to obtain more useful sparsity at an acceptable
quality cost? Answering this question matters when choosing which components
of a recipe to retain, where to invest additional training, and which
operations warrant sparse-kernel development.

The available approaches act on different parts of this problem. Post-hoc
thresholding exploits the tolerance of an existing representation;
activation regularization changes what the model learns; and broader site
placement exposes additional operations to sparsification. TEAL, ProSparse,
and Q-Sparse establish the effectiveness of these approaches and investigate
choices within them~\citep{liu2024teal,song2024prosparse,wang2024qsparse}.
These findings motivate examining the additional value of pressure as
thresholding and placement change.
An intervention that helps on its own may add little to another intervention,
or impose a different quality cost when applied elsewhere. Without knowing
these dependencies, an improved sparsity fraction gives limited guidance
about why a recipe works or how to adapt it to a different model.

We study this problem through three separately controllable design choices:
\emph{pressure}, which encourages selected activations toward zero;
\emph{thresholding}, which sets values to exactly zero according to a gate;
and \emph{site placement}, which specifies where each acts.
\emph{The interventions are separately controllable, but their effects need
not be separable.} Our central question is how their individual effects and
interactions determine the quality--sparsity trade-off. In particular, does
adding pressure become more or less useful as thresholding changes, and how
does that response depend on where pressure and gates are applied?

Explaining that response requires connecting learned behavior to the
computation it affects. A reduction in activation magnitude can leave the
number of exact zeros unchanged; additional zeros can also be concentrated
at sites that affect only a small share of the forward pass. Conversely,
extending the gate set can expose additional operations to sparsification while
changing the representation the model must learn. The empirical issue is how
learning responds to these choices: when pressure produces additional zeros,
how much quality that change costs, and which operations account for the
result. Distinguishing these outcomes helps diagnose whether limited gains
reflect the response to pressure, the computational reach of the chosen sites,
or the cost of expanding the intervention.

To make that diagnosis model-wide, we pair observed zero-product counts with
an architecture-dependent reach ceiling. $\Smodel$ measures the
fraction of counted forward matrix-multiplication products with a zero
activation operand; $\Smax$ describes the fraction
reachable if the selected sites were entirely zero. We refer to the former as
\emph{model-wide sparsity}, distinguishing it from a local activation zero
fraction. This accounting makes the location of sparsity consequential: the
same local zero rate can affect different amounts of computation as the
architecture or workload changes. It helps distinguish a change in learned
sparsity from a change in the work reachable through the chosen sites.
Direct timing then tests
whether an implementation can exploit the observed zeros on hardware.

We organize matched pretraining comparisons around a \emph{sparsification
intervention ladder} (Figure~\ref{fig:intervention-ladder}), which organizes gate
and pressure alternatives across feed-forward (FFN) and attention sites. Pythia-family
models provide a common architecture for examining these choices and their
behavior across sizes~\citep{biderman2023pythia}. We use the ladder to trace
individual and joint intervention effects from learned activations to the
operations they affect. A transferable explanation must account for changes
in those effects at another model size, where both learned responses and
computational reach can differ. An account that meets this requirement can
guide how sparsification recipes are adapted, including when to strengthen
pressure, change a gate, or extend placement.
